\section{Audit against all existing proposal bundles}
\label{sec:exclusion}

The following matrix records the operative mechanism already proposed in
each directory. The archive includes the fuller audit and a manifest of
59 relevant text files with their Git blob identifiers and SHA-256 hashes.
Future-direction sections were also checked; a merely restated future
direction was not counted as a new proposal.

\small
\begin{longtable}{@{}p{.35\textwidth}p{.59\textwidth}@{}}
\toprule
Existing bundle & Mechanism excluded from novelty here\\
\midrule
\endfirsthead
\toprule Existing bundle & Mechanism excluded from novelty here\\ \midrule
\endhead
\proposal{centered_theta_lift_simplification}{Centered theta lifts} &
Centered covariance, direct selection in the ramified source, adaptive
completion, and random contraction products.\\
\proposal{cm_norm_theta_simplification}{CM norm theta} &
Two/three-factor local cancellation and four-copy sharpness, with padding
to the old interface.\\
\proposal{cyclotomic-interpolation-and-theta-lifts}{Cyclotomic interpolation} &
Dimension-independent multivariable torsion interpolation and improved
near-spanning.\\
\proposal{discrete_hecke_simplification}{Discrete Hecke} &
Integer Pieri-cone enumeration, Smith saturation, and possible direct
counting in place of Frobenius packing.\\
\proposal{edge_hecke_simplification}{Edge Hecke} &
Orbit-trace invariant averages and principal-submatrix bounds for
exceptional Smith directions.\\
\proposal{gram_adjoint_simplification}{Gram adjoint} &
Integral Fock Gram adjoints and removal of specialized tight-frame
bookkeeping in one-lift estimates.\\
\proposal{grh-class-count-finite-quotient}{Finite-quotient class count} &
Class encoding by short-basis Gram data and residues in a finite quotient.\\
\proposal{grh_order_tests_2026-11-27}{Order tests (27 November)} &
Norm-divisibility order tests and subgroup exclusions instead of primitive
finite-field elements.\\
\proposal{newton_torsion_simplification}{Newton torsion} &
One-shell divided differences, prefix denominator budgets, and a singly
exponential order ceiling.\\
\proposal{orbital_theta_exchange}{Orbital theta exchange} &
Canonical orbital sections, algebraic twisted coinvariants, and
complementary-frame matching.\\
\proposal{proof_simplifications}{Proof simplifications} &
Quadratic twist averages, direct Perron estimates, and ramified-source
spanning.\\
\proposal{rh-simplifications}{RH simplifications} &
Sharp rank recurrences, torsion interpolation, and large-sieve twist
selection.\\
\proposal{rh_sparse_theta_proof}{Sparse theta proof} &
Sparse-cutoff interpolation and Gaussian second-moment twist selection.\\
\proposal{rh_two_order_tests_proposal}{Two-order tests} &
Weaker finite-field label tests and a weakened common-scalar-order
interface.\\
\proposal{robust_theta_labels}{Robust theta labels} &
Three power tests, two-prime scalar certificates, and stronger label counts.\\
\proposal{shorter-theta-lift-proof}{Shorter theta-lift proof} &
Elementary or large-sieve twist means, direct spanning, and far-cutoff
prime estimates.\\
\proposal{signed_theta_gain_simplification}{Signed theta gain} &
Signed local gains, avoiding rounded rational short-form normalization,
and product-formula assembly with improved colength budgets.\\
\proposal{theta-lift-simplifications}{Theta-lift simplifications} &
Flag-adapted integral rows, compressed determinants, denser labels, and
explicit sparse infinite places.\\
\bottomrule
\end{longtable}
\normalsize

The nearest overlaps deserve a precise distinction:
\begin{itemize}
\item The normalization in signed theta gains is the rational divisor
used to turn the short forms into short integral vectors. The present
proposal removes exact modulus metadata of exchange scalars in the
lower-volume branch.
\item Orbital theta exchange changes how the square laws are constructed.
The present cost calculus consumes whatever exact laws that construction
provides. Better local costs can be passed through unchanged.
\item The existing Gram-adjoint proposal and the one-vector branching
argument concern different operators and different source modules.
\item A rank--height budget, an exterior-power replacement for completion,
and stronger torsion interpolation are already explicit future directions;
they are excluded from the present proposals.
\end{itemize}

\section{Implementation order and acceptance criteria}

\subsection{A staged migration}

\begin{enumerate}
\item Add exact power homogeneity next to the existing
\code{HeightLe.pow}. Compile the specialization in a fully built
checkout and retain the current height definition.
\item Add the guarded squared-cost interface. Prove constants, products,
one-crossing costs, and reciprocal zero-cost lemmas for the specific
norm-power factors. During this stage, the old exact-modulus lemmas can
serve as adapters, making the new consumer interface easy to validate.
\item Change the triangle's scalar grade to additive costs, retaining its
mixed Gram statement, injectivity proof, and exact weights. Add the
determinant recurrence rather than a public accumulated Rallis product.
\item Convert the chain and auxiliary-form consumers. Preserve the
boundary Euler-factor cancellation, \eqref{eq:oldbound}, and the exact
statement of \code{integral_forms_long}.
\item Once no consumers need the exact ratio metadata, delete it and
measure net source size. Do not count a wrapper around the old development
as a completed simplification.
\item Independently replace the branching norm proof by the
Gram-adjoint/Schur argument, calibrated by the existing Cartan norm
inequality or by the direct fixed-tensor lemma.
Initially preserve the old public theorem signature; drop the redundant
weight hypothesis in a later cleanup if useful.
\end{enumerate}

If a small, self-contained mathematical replacement is the immediate
priority, implement Proposal II first. If the priority is reducing
coupling across the determinant and auxiliary-form developments, implement
Proposal I first. The two changes do not depend on each other.

\subsection{Concrete verification gates}

The acceptance criterion for Proposal I is the current Theorem C with
the same \(H-3L\), together with the current final five statements.
Check the zero-target branch, actual scalar nonvanishing, the inverse
turn scalar, common twist quantifiers, and the conductor boundary
cancellation explicitly. Those are the places where a seemingly
shorter proof can silently change its meaning.

For Proposal II, the acceptance criterion is the current squared norm
bound in \code{IntegralMap}, with coefficient \(c^2\) and no new
dimension or factorial factor. Check that the calibration vectors
(or the alternative fixed tensor) have norm exactly one and that the
target commutant theorem is applied to the correctly twisted action.
A loose bound in the calibration would defeat the proposed simplification.

After both integrations, build the affected closure, build the full project,
print the axioms of all five public theorems, and replay the desired
independent checker workflow. The archive's verification script checks
only the new standalone components. It deliberately does not label
those checks a full GRH verification.

\section{Tempting alternatives that fail this audit}

Several apparently larger shortcuts were investigated and rejected.
Recording the obstruction is useful because these alternatives can look
more dramatic than the actual proposals.

\paragraph{Multiply ordinary Rallis formulas along the chain.}
The stable-rank hypothesis does not hold for every consecutive lift.
The triangular exchange and adjoint identities supply missing information,
and their nonzero scalar analysis is not optional.

\paragraph{Use a full determinant lower bound on an almost basis.}
A full determinant controls a product of all singular values.
An arbitrary omitted complement can carry compensating large singular
values. The auxiliary completion in the source supplies exactly the
missing control. Removing it requires a new exterior-power theorem,
a direction already mentioned among the existing proposals.

\paragraph{Replace the Fitting calculation by a square-zero quotient.}
The source already uses the ideal generated by the squares of the
deformation variables in its
\src{New/Congruences/FirstOrder/FittingBlock}{\code{FittingBlock}} development.
The elementary annihilator calculation in that quotient can improve
presentation, but it does not remove the difficult lifting or Fitting
inputs. It is not a substantial new simplification of this proof.

\paragraph{Recover integral characteristic data from power traces for free.}
Newton identities introduce denominators. For instance, over
\(\Z_p[t]\), the matrix \(tI_p\) has traces \(pt^k\) and determinant \(t^p\).
Modulo \(p\), every nonconstant generator \(pt^k\) vanishes, whereas
\(t^p\) does not. Thus \(t^p\) is not in the
\(\Z_p\)-algebra generated by those traces. A factorial paid to clear
denominators also scales the distinguished deformation coefficient;
it is not free precision in the original valuation budget.
The archive contains a fuller note on this obstruction.

\section{Delivered artifacts and reproducibility}

The archive contains:
\begin{itemize}
\item this report as \code{grh_simplification.tex}, its included section
sources, and \code{grh_simplification.pdf};
\item the two checked Lean files
\code{lean/HeightPower.lean} and \code{lean/MatrixChainBudget.lean};
\item the separately labeled source-integration candidate for
\code{New.HeightLe};
\item build instructions, pinned toolchain/dependency configuration,
and scripts for compiling the report and checking the standalone Lean;
\item source/proposal manifests, source metrics, verification logs,
and the mathematical and exclusion-review notes;
\item a SHA-256 manifest for the delivered files.
\end{itemize}

The source repository is not included in the archive. Immutable source
links and hashes identify the material analyzed. The proof files and
source-integration candidate contain no placeholders standing in for
the missing full migration; that migration is described as future work.

The smallest immediate addition is \code{HeightLe.of_pow}: it is exact,
independently useful, and already checked in a faithful generalization.
Its payoff comes when it enables the determinant branch to retain the
height information it actually consumes. The branching argument offers
a separate local mathematical reduction whose proof is shorter because
equivariance makes an all-vector calculation unnecessary.
